\documentclass[12pt,a4paper]{article}
\usepackage[utf8]{inputenc}
\usepackage[T1]{fontenc}
\usepackage[french]{babel}
\usepackage{geometry}
\usepackage{xcolor}
\usepackage{fancyhdr}

\geometry{margin=2cm}
\pagestyle{fancy}
\fancyhf{}
\fancyhead[L]{CEJM - Fiche de Révision}
\fancyhead[R]{Chapitre 3}
\fancyfoot[C]{\thepage}
\setlength{\headheight}{15pt}

% Couleurs
\definecolor{bleu}{RGB}{0,90,170}
\definecolor{vert}{RGB}{0,150,0}
\definecolor{rouge}{RGB}{200,0,0}
\definecolor{gris}{RGB}{240,240,240}

% Commandes pour encadrés
\newcommand{\definition}[2]{\noindent\fcolorbox{bleu}{gris}{\parbox{\dimexpr\textwidth-2\fboxsep-2\fboxrule\relax}{\textbf{\textcolor{bleu}{DÉFINITION :}} #1\\\textit{#2}}}}

\newcommand{\pointcle}[1]{\noindent\fcolorbox{vert}{white}{\parbox{\dimexpr\textwidth-2\fboxsep-2\fboxrule\relax}{\textbf{\textcolor{vert}{POINT CLÉ :}} #1}}}

\newcommand{\exemple}[1]{\noindent\fcolorbox{rouge}{white}{\parbox{\dimexpr\textwidth-2\fboxsep-2\fboxrule\relax}{\textbf{\textcolor{rouge}{EXEMPLE :}} #1}}}

\title{\textbf{Fiche de Révision CEJM}\\[2mm]\large Chapitre 3 : La formation du contrat}
\author{}
\date{}

\begin{document}
\maketitle

\section*{Problématique}
\textbf{Comment se forme juridiquement un contrat et quelles sont les conditions de sa validité ?}

\section{Définitions essentielles}

\definition{Contrat}{Accord de volontés entre deux ou plusieurs personnes destiné à créer des obligations juridiques.}

\vspace{3mm}
\definition{Offre (pollicitation)}{Proposition ferme et précise de contracter, faite à une personne déterminée ou au public.}

\vspace{3mm}
\definition{Acceptation}{Adhésion pure et simple aux termes de l'offre, exprimée de manière expresse ou tacite.}

\vspace{3mm}
\definition{Consentement}{Accord de volonté libre et éclairé des parties au contrat.}

\vspace{3mm}
\definition{Cause}{Raison, motif qui pousse à contracter. Elle doit être licite et morale.}

\section{Les conditions de formation du contrat}

\subsection{Conditions de fond (validité)}
\pointcle{4 conditions essentielles pour la validité d'un contrat}

\subsubsection{1. Le consentement}
\begin{itemize}
    \item \textbf{Libre} : sans contrainte physique ou morale
    \item \textbf{Éclairé} : en connaissance de cause
    \item \textbf{Exprès ou tacite} : par parole, écrit ou comportement
\end{itemize}

\textbf{Vices du consentement :}
\begin{itemize}
    \item \textbf{Erreur} : méconnaissance de la réalité
    \item \textbf{Dol} : tromperie volontaire d'une partie
    \item \textbf{Violence} : contrainte physique ou morale
\end{itemize}

\subsubsection{2. La capacité}
\begin{itemize}
    \item \textbf{Capacité de jouissance} : aptitude à avoir des droits
    \item \textbf{Capacité d'exercice} : aptitude à exercer ses droits
    \item \textbf{Incapacités} : mineurs, majeurs protégés
\end{itemize}

\subsubsection{3. L'objet}
\begin{itemize}
    \item Doit être \textbf{déterminé ou déterminable}
    \item Doit être \textbf{possible} (physiquement et juridiquement)
    \item Doit être \textbf{licite} (conforme à la loi et aux bonnes mœurs)
\end{itemize}

\subsubsection{4. La cause}
\begin{itemize}
    \item \textbf{Existence} : motif réel de l'engagement
    \item \textbf{Licéité} : conforme à la loi
    \item \textbf{Moralité} : conforme aux bonnes mœurs
\end{itemize}

\subsection{Conditions de forme}
\pointcle{Généralement, le contrat est consensuel (pas de forme particulière)}
\begin{itemize}
    \item \textbf{Contrats consensuels} : accord de volonté suffit
    \item \textbf{Contrats solennels} : forme particulière exigée (notaire, écrit...)
    \item \textbf{Contrats réels} : remise de la chose nécessaire
\end{itemize}

\section{Le processus de formation}

\subsection{La phase de négociation}
\pointcle{Pourparlers précontractuels sans engagement}
\begin{itemize}
    \item Liberté de négocier
    \item Obligation de bonne foi
    \item Responsabilité en cas de rupture abusive
\end{itemize}

\subsection{L'offre de contracter}
\pointcle{Proposition ferme, précise et complète}
\begin{itemize}
    \item \textbf{Ferme} : intention réelle de s'engager
    \item \textbf{Précise} : éléments essentiels déterminés
    \item \textbf{Complète} : tous les éléments du futur contrat
\end{itemize}

\textbf{Durée de l'offre :}
\begin{itemize}
    \item Offre avec délai : maintenue pendant le délai
    \item Offre sans délai : délai raisonnable
    \item Révocation possible avant acceptation
\end{itemize}

\subsection{L'acceptation}
\pointcle{Adhésion pure et simple aux termes de l'offre}
\begin{itemize}
    \item \textbf{Pure et simple} : sans modification
    \item \textbf{Expresse} : orale ou écrite
    \item \textbf{Tacite} : par comportement non équivoque
\end{itemize}

\section{Le moment et le lieu de formation}

\subsection{Moment de formation}
\pointcle{Le contrat est formé dès la rencontre de l'offre et de l'acceptation}
\begin{itemize}
    \item \textbf{Entre présents} : immédiatement
    \item \textbf{Entre absents} : réception de l'acceptation
    \item \textbf{Électronique} : accusé de réception
\end{itemize}

\subsection{Lieu de formation}
\begin{itemize}
    \item Lieu où l'acceptation a été donnée
    \item Important pour déterminer la juridiction compétente
\end{itemize}

\section{Types de contrats}

\subsection{Classification selon les obligations}
\begin{itemize}
    \item \textbf{Contrat synallagmatique} : obligations réciproques
    \item \textbf{Contrat unilatéral} : obligations d'une seule partie
\end{itemize}

\subsection{Classification selon l'intérêt}
\begin{itemize}
    \item \textbf{Contrat à titre onéreux} : avantages réciproques
    \item \textbf{Contrat à titre gratuit} : libéralité d'une partie
\end{itemize}

\subsection{Classification selon la formation}
\begin{itemize}
    \item \textbf{Contrat de gré à gré} : négociation libre
    \item \textbf{Contrat d'adhésion} : conditions imposées
\end{itemize}

\section{Exemples concrets}

\exemple{Vente en ligne : offre sur le site web, acceptation par commande, formation à la validation}

\exemple{Contrat de travail : négociation des conditions, signature du contrat écrit obligatoire}

\exemple{Achat en magasin : offre par l'étiquetage, acceptation par le passage en caisse}

\section{Sanctions en cas de défaut}

\subsection{Nullité du contrat}
\pointcle{Anéantissement rétroactif du contrat}
\begin{itemize}
    \item \textbf{Nullité absolue} : ordre public (5 ans)
    \item \textbf{Nullité relative} : protection d'une partie (5 ans)
\end{itemize}

\subsection{Conséquences}
\begin{itemize}
    \item Restitutions réciproques
    \item Dommages-intérêts possibles
    \item Confirmation possible (nullité relative)
\end{itemize}

\section{Points clés à retenir}

\pointcle{Un contrat valide nécessite : consentement, capacité, objet, cause}

\pointcle{Le consentement doit être libre et éclairé}

\pointcle{L'offre doit être ferme, précise et complète}

\pointcle{L'acceptation doit être pure et simple}

\pointcle{Le contrat se forme à la rencontre de l'offre et de l'acceptation}

\section{Mots-clés à connaître}
\begin{itemize}
    \item Contrat/Convention
    \item Offre/Pollicitation
    \item Acceptation
    \item Consentement
    \item Capacité juridique
    \item Objet du contrat
    \item Cause du contrat
    \item Vices du consentement
    \item Nullité absolue/relative
    \item Contrat consensuel/solennel
\end{itemize}

\end{document}